\section{Indução e Recursão transfinitas}\label{sec:inducaoERecursaoTransfinitas}
\index{indução transfinita}
\index{recursão transfinita}
Vimos que $\in$ é uma boa ordem sobre a classe dos ordinais.
Em particular, vimos que se $X$ é um conjunto de ordinais, então $X$ possui um elemento mínimo.

\begin{proposition}[Esquema de boa ordenação]\label{prop:esquemaDaBoaOrdenacao}\index{esquema de boa ordenação}
    Seja $\phi(\alpha, t_1, \dots, t_n)$ uma fórmula.
    Suponha que, para alguns $t_1, \dots, t_n$ fixados, exista um ordinal $\alpha$ tal que $\phi(\alpha, t_1, \dots, t_n)$.

    Então existe o menor ordinal $\alpha$ tal que $\phi(\alpha, t_1, \dots, t_n)$.

    Em símbolos:

    \[
        \forall t_1 \dots \forall t_n \Big( \exists \alpha \in \ON \phi(\alpha, t_1, \dots, t_n) \rightarrow \exists! \alpha \in \ON (\phi(\alpha, t_1, \dots, t_n) \wedge \forall \beta <\alpha (\neg\phi(\beta, t_1, \dots, t_n)))\Big).
    \]
\end{proposition}
\begin{proof}
    Para a existência, seja $\xi$ um ordinal tal que $\phi(\xi, t_1, \dots, t_n)$.
    Se $\forall \alpha <\xi \neg\phi(\alpha, t_1, \dots, t_n)$, então está provada a existência.
    Se não, considere $X=\{\alpha <\xi : \phi(\alpha, t_1, \dots, t_n)\}$.
    Temos que $X$ é um conjunto de ordinais não vazio, logo, possui elemento mínimo $\alpha$.
    Segue que para todo $\beta<\alpha$, temos $\beta\notin X$, mas $\beta<\xi$, de modo que $\neg\phi(\beta, t_1, \dots, t_n)$.
    Assim, $\alpha$ é como queríamos.

    Para a unicidade, suponha que $\alpha'$ também satisfaça
    \[
        \phi(\alpha', t_1,\dots,t_n)\land \forall \beta<\alpha'\,\neg\phi(\beta,t_1,\dots,t_n).
    \]
    Se $\alpha'<\alpha$, então, como $\phi(\alpha', t_1,\dots,t_n)$, isso contradiz a minimalidade de $\alpha$.
    Analogamente, não pode ocorrer $\alpha<\alpha'$.
    Logo, $\alpha=\alpha'$.
\end{proof}

A indução transfinita nada mais é do que a aplicação sistemática do princípio do menor contraexemplo aos ordinais. A ideia é simples: para provar que uma propriedade vale para todo ordinal, basta mostrar que ela vale em um ordinal sempre que já valha em todos os ordinais anteriores.

Assim como na indução usual sobre $\omega$, o que se faz é excluir a possibilidade de um primeiro contraexemplo. A diferença é que agora o domínio de indução não é apenas $\omega$, mas toda a classe dos ordinais.

\begin{proposition}[Esquema da indução transfinita]\label{prop:esquemaDaInducaoTransfinita}
    Seja $\phi(\alpha, t_1, \dots, t_n)$ uma fórmula.
    
    Suponha que, para alguns $t_1, \dots, t_n$ fixados e para todo ordinal $\alpha$, se $\forall \beta <\alpha \phi(\beta, t_1, \dots, t_n)$, então $\phi(\alpha, t_1, \dots, t_n)$.

    Então $\forall \alpha \in \ON \phi(\alpha, t_1, \dots, t_n)$.
    
    Em símbolos:
    \[
        \forall t_1 \dots \forall t_n \Big( \forall \alpha \in \ON ((\forall \beta <\alpha \phi(\beta, t_1, \dots, t_n)) \rightarrow \phi(\alpha, t_1, \dots, t_n)) \rightarrow \forall \alpha \in \ON \phi(\alpha, t_1, \dots, t_n)\Big).
    \]
\end{proposition}
\begin{proof}
    Fixe $t_1, \dots, t_n$ e suponha que a tese é falsa.
    Então existe um ordinal $\alpha$ tal que $\neg\phi(\alpha, t_1, \dots, t_n)$.
    Pela Proposição~\ref{prop:esquemaDaBoaOrdenacao}, existe o menor ordinal $\alpha$ tal que $\neg\phi(\alpha, t_1, \dots, t_n)$.
    Assim, para todo $\beta <\alpha$, temos $\phi(\beta, t_1, \dots, t_n)$.
    Mas então $\phi(\alpha, t_1, \dots, t_n)$, o que é uma contradição.
\end{proof}

\begin{definition}\label{def:sequenciaTransfinita}\index{sequência transfinita}
    Uma sequência transfinita é uma função cujo domínio é um ordinal.
    Se $s$ é uma sequência transfinita, escrevemos $s \in V^{<\ON}$.
\end{definition}

É importante distinguir indução transfinita de recursão transfinita.

A indução transfinita é um \emph{método de demonstração}: ela serve para provar que uma propriedade vale para todos os ordinais.

Já a recursão transfinita é um \emph{método de definição}: ela serve para construir uma função $F$ definida em todos os ordinais, de modo que o valor $F(\alpha)$ seja determinado a partir dos valores anteriores, $F|\alpha$.

Um importante exemplo, sobre o qual discutiremos mais a fundo na próxima seção, é a \emph{soma de ordinais}, que é uma extensão natural da soma usual de números naturais.
Ela é definida a partir das seguintes propriedades, em que $\alpha$ e $\beta$ são ordinais:
\begin{itemize}
    \item $\beta+0=\beta$;
    \item $\beta+S(\alpha)=S(\beta+\alpha)$;
    \item $\beta+\alpha=\sup\{\beta+\xi : \xi <\alpha\}$, se $\alpha$ é um ordinal limite.
\end{itemize}
Ou seja, começa-se definindo $\beta+0=\beta$.
Definido $\beta+\alpha$, define-se $\beta+S(\alpha)$ como o ordinal sucessor de $\beta+\alpha$.
Finalmente, se $\alpha$ é um ordinal limite, define-se $\beta+\alpha$ como o supremo de $\{\beta+\xi : \xi <\alpha\}$.

Para se executar essa definição, para se definir $\beta+\alpha$, é necessário ter o conhecimento prévio do valor de potencialmente todos os valores $\beta+\xi$ para $\xi <\alpha$.

Diferentemente do que acontece na recursão de números naturais, o objeto resultante não é uma função.
Se fosse, no caso, a função $+$ teria domínio $\ON\times \ON$, que não é um conjunto.
O objeto resultante é um símbolo de função $+$ introduzido por definição, ou, analogamente, uma definição para introduzir tal símbolo.

\begin{metatheorem}[Esquema da recursão transfinita]\label{mthm:esquemaDaRecursaoTransfinita}\index{recursão transfinita}
    Suponha que um símbolo de função binário $G(p, s)$ seja introduzido por definição de modo que
    \[
        \forall s \in V^{<\ON}\forall p\forall y
        \Bigl(G(p,s)=y \leftrightarrow \phi(s,y,p)\Bigr).
    \]

    Então é possível introduzir um símbolo de função binário $F(p, \alpha)$ de modo que
    \[
        \forall \alpha \in \ON \forall p \Big( F_{p}(\alpha)=G(p, F_p|\alpha)\Big).
    \]

    Ademais, se $H$ é um símbolo de função binário adicionado por definição tal que
    \[
        \forall p \forall \alpha \in \ON \Bigl( H_p(\alpha)=G(p, H_p|\alpha)\Bigr),
    \]
    então
    \[
        \forall p \forall \alpha \in \ON \Bigl( F_p(\alpha)=H_p(\alpha)\Bigr).
    \]
\end{metatheorem}
Chamamos o resultado dado pelo Metateorema~\ref{mthm:esquemaDaRecursaoTransfinita} de metateorema porque ele não afirma apenas uma propriedade de conjuntos já dados na teoria: ele afirma que, a partir de um símbolo de função $G$ introduzido por definição, é possível introduzir um novo símbolo de função $F$ satisfazendo a equação recursiva desejada.

Ou seja, o conteúdo do resultado é que o processo de definição recursiva transfinita pode ser formalizado dentro da teoria. A partir disso, os resultados concretos sobre a função assim definida passam a ser teoremas internos da teoria.

Antes de apresentar uma demonstração do metateorema, vejamos um exemplo de aplicação do mesmo.
\begin{example}\label{exa:somaOrdinal}
    Considere o operador binário $G$ definido por:
    \[
        G(p, s)=\begin{cases}
            p & \text{se } s=\emptyset,\\
            S(s(\alpha)) & \text{se } s\text{ é função } \land  \exists \alpha \in \ON \dom(s)=\alpha+1,\\
            \bigcup_{\xi \in \alpha} s(\xi) & \text{se } s\text{ é função } \land \exists \alpha \in \ON \dom(s)=\alpha \text{ é limite}.
        \end{cases}
    \]

    O esquema da recursão nos dá uma receita para introduzir por definição $F(p, \alpha)$ tal que $F_p(\alpha)=G(p, F_p|\alpha)$ para todo $p$ e $\alpha$.
    Se $\beta$ é um ordinal, denotando $F_\beta(\alpha)=\beta+\alpha$, temos que:

    \begin{itemize}
        \item $\beta+0=G(\beta, F_\beta|0)=\beta$;
        \item $\beta+S(\alpha)=S(\beta+\alpha)$;
        \item $\beta+\alpha=\sup\{\beta+\xi : \xi <\alpha\}$, se $\alpha$ é um ordinal limite.
    \end{itemize}

    Por uma indução transfinita simples, vê-se que $\beta+\alpha \in \ON$ para todo $\beta, \alpha \in \ON$: suponha que para todo ordinal $\xi<\alpha$, $\beta+\xi$ é um ordinal.
    O mesmo vale para $\alpha$. Com efeito, se $\alpha=0$, temos que $\beta+0=\beta$ é um ordinal.
    Se $\alpha=\xi+1$ é um ordinal sucessor, por hipótese de indução, $\beta+\xi$ é um ordinal, logo, $S(\beta+\xi)$ é um ordinal.
    Finalmente, se $\alpha$ é um ordinal limite, por hipótese de indução, $\beta+\xi$ é um ordinal para todo $\xi <\alpha$, logo, $\beta+\alpha=\bigcup\{\beta+\xi : \xi <\alpha\}$ é um ordinal.
    
    Portanto, $\beta+\alpha=\sup\{\beta+\xi : \xi <\alpha\}$ se $\alpha$ é um ordinal limite.
\end{example}

Dedicaremos o resto da seção para a prova do Metateorema~\ref{mthm:esquemaDaRecursaoTransfinita}.
Iniciaremos pelo seguinte esquema de definições.
\begin{definition}\label{def:computacaoParcialTransfinita}
    Seja $G(p, s)$ um símbolo de função binário introduzido por definição de modo que
    \[
        \forall s \in V^{<\ON}\forall p\forall y
        \Bigl(G(p,s)=y \leftrightarrow \phi(s,y,p)\Bigr).
    \]

    Para todo ordinal $\alpha$ e todo $p$, uma computação parcial de comprimento $\alpha$ e parâmetro $p$ baseada em $G$ é uma função $h$ tal que $\dom(h)=\alpha$ e para todo $\beta <\alpha$, $h(\beta)=G(p, h|\beta)$.
    
    Nesse caso, escrevemos $\Ap_G(h, \alpha, p)$.
    Em símbolos:

    \[
        \forall h\forall \alpha\forall p
        \Bigl(\Ap_G(h,\alpha,p) \leftrightarrow
        h \text{ é função} \land \alpha\in\ON \land \dom(h)=\alpha \land
        \forall \beta<\alpha\bigl(h(\beta)=G(p,h|\beta)\bigr)\Bigr).
    \]
\end{definition}
Uma computação parcial de comprimento $\alpha$ deve ser vista como uma tentativa de construir a função recursiva apenas até o estágio $\alpha$.
Ela contém, para cada $\beta<\alpha$, exatamente o valor que a regra recursiva prescreve a partir do trecho já construído $h|\beta$.


Agora provaremos o seguinte esquema de lemas.
\begin{lemma}\label{lem:compatibilidadeDeComputacoesParciais}
    Seja $G(p, s)$ um símbolo de função introduzido por definição.
    Para todos $p, \alpha, \alpha', h, h'$, se $\Ap_G(h, \alpha, p)$ e $\Ap_G(h', \alpha', p)$, então $h$ e $h'$ são compatíveis.
\end{lemma}
\begin{proof}
    Fixe $p, \alpha, \alpha', h, h'$ tais que $\Ap_G(h, \alpha, p)$ e $\Ap_G(h', \alpha', p)$.

    Se $h$ e $h'$ não são compatíveis, então existe $\beta\in \alpha\cap \alpha'=\min\{\alpha, \alpha'\}$ tal que $h(\beta)\neq h'(\beta)$.
    Tome $\beta$ o menor tal ordinal.

    Então $\beta\subseteq \alpha\cap \alpha'$ e $h|\beta=h'|\beta$, de modo que
    \[
        h(\beta)=G(p, h|\beta)=G(p, h'|\beta)=h'(\beta),
    \]
    o que é uma contradição.
\end{proof}

\begin{corollary}\label{cor:unicidadeDeComputacaoParcial}
    Seja $G(p, s)$ um símbolo de função introduzido por definição.
    Para todos $p, \alpha, h, h'$, se $\Ap_G(h, \alpha, p)$ e $\Ap_G(h', \alpha, p)$, então $h=h'$.
\end{corollary}
\begin{proof}
    Se $\Ap_G(h, \alpha, p)$ e $\Ap_G(h', \alpha, p)$, então $h$ e $h'$ são compatíveis.
    Mas $\dom(h)=\dom(h')=\alpha$, logo $h=h'$.
\end{proof}
Resta agora mostrar que tais computações parciais de fato existem para todo comprimento ordinal.
Uma vez estabelecida essa existência, poderemos definir o valor da iteração no estágio $\alpha$ escolhendo qualquer computação parcial suficientemente longa e lendo nela o valor na posição $\alpha$.

A compatibilidade provada no Lema~\ref{lem:compatibilidadeDeComputacoesParciais} garantirá que essa escolha independe da computação parcial utilizada.
\begin{lemma}\label{lem:existenciaDeComputacoesParciais}
    Seja $G(p, s)$ um símbolo de função introduzido por definição.
    Para todo $p$ e todo ordinal $\alpha$, existe uma computação parcial de comprimento $\alpha$ e parâmetro $p$ baseada em $G$.
\end{lemma}
\begin{proof}
    Fixe $p$.
    Procederemos por indução transfinita sobre $\alpha$.

    Suponha que para todo $\beta <\alpha$, existe uma computação parcial de comprimento $\beta$ e parâmetro $p$ baseada em $G$.
    Vejamos que o mesmo vale para $\alpha$.

    Se $\alpha=0$, tome $h=\emptyset$.

    Se $\alpha=\beta+1$ é um ordinal sucessor, por hipótese de indução, existe uma computação parcial $h$ de comprimento $\beta$ e parâmetro $p$ baseada em $G$.
    Então, tome $h'=h\cup\{(\beta, G(p, h))\}$.
    Assim, $h$ é função, $\dom(h')=\beta\cup\{\beta\}=\alpha$ e para todo $\gamma <\beta$, $h'(\gamma)=h(\gamma)=G(p, h|\gamma)=G(p, h'|\gamma)$.
    Finalmente, $h'(\beta)=G(p, h)=G(p, h'|\beta)$.
    Assim, $\Ap_G(h', \alpha, p)$.

    Finalmente, se $\alpha$ é um ordinal limite, para cada $\beta <\alpha$, por hipótese de indução e pelo Corolário~\ref{cor:unicidadeDeComputacaoParcial}, existe uma única computação parcial $h_\beta$ de comprimento $\beta$ e parâmetro $p$ baseada em $G$.
    Pelo Axioma~\ref{ax:substituicao}, seja $H=\{h_\beta : \beta <\alpha\}$.
    
    Pelo Lema~\ref{lem:compatibilidadeDeComputacoesParciais}, $H$ é um conjunto de funções compatíveis, logo, $h=\bigcup H$ é uma função de domínio $\bigcup_{\beta <\alpha} \beta=\alpha$.
    Finalmente, se $\beta <\alpha$, então $\beta+1\leq \alpha$.
    Como $\alpha$ é limite, temos que $\beta+1<\alpha$.
    Temos que $h(\beta)=h_{\beta+1}(\beta)=G(p, h_{\beta+1}|\beta)=G(p, h|\beta)$.
    Assim, $\Ap_G(h, \alpha, p)$.
\end{proof}

\begin{definition}\label{def:iteracaoTransfinita}
    Seja $G(p, s)$ um símbolo de função introduzido por definição.
    A $\alpha$-ésima iteração de $G$ com parâmetro $p$, denotada por $F^G(p, \alpha)$ (ou simplesmente $F(p, \alpha)$ se $G$ for claro do contexto), é $h(\alpha)$, em que $h$ é alguma computação parcial de comprimento $\beta>\alpha$ e parâmetro $p$ baseada em $G$.
    Em símbolos:
    \[
        \forall p\forall \alpha\in\ON\forall y
        \Bigl(F^G(p,\alpha)=y \leftrightarrow
        \exists h\exists \beta\bigl(\Ap_G(h,\beta,p)\land \alpha<\beta \land h(\alpha)=y\bigr)\Bigr).
    \]
\end{definition}
Pelos Lemas~\ref{lem:compatibilidadeDeComputacoesParciais} e \ref{lem:existenciaDeComputacoesParciais}, $F^G(p, \alpha)$ é bem definido, pois tal $y$ existe e é único.
Assim, temos o seguinte metateorema.
\begin{theorem}\label{thm:equacaoDaIteracaoTransfinita}
    Seja $G(p, s)$ um símbolo de função binário introduzido por definição.

    Denote por $F_p(\alpha)=F(p, \alpha)$ a $\alpha$-ésima iteração de $G$.

    Então
    \[
        \forall p \forall \alpha \in \ON \Bigl( F_p(\alpha)=G(p, F_p|\alpha)\Bigr).
    \]
\end{theorem}
\begin{proof}
    Seja $p$ e $\alpha$ fixados.
    Seja $h$ uma computação parcial de comprimento $\beta>\alpha$ e parâmetro $p$ baseada em $G$ tal que $h(\alpha)=F_p(\alpha)$.
    Pela definição de $F$, temos que $h(\xi)=F_p(\xi)$ para todo $\xi <\alpha$.
    Assim, $h|\alpha=F_p|\alpha$.

    Desta forma, $F_p(\alpha)=h(\alpha)=G(p, h|\alpha)=G(p, F_p|\alpha)$.
\end{proof}
O passo final consiste em mostrar que a solução recursiva obtida é única.
Isto é, qualquer outra função que satisfaça a mesma equação recursiva deve coincidir com $F$ em todos os ordinais, independente de como ela seja definida.

Novamente, a ferramenta apropriada para isso é a indução transfinita.
\begin{theorem}\label{thm:unicidadeDaIteracaoTransfinita}
    Seja $G(p, s)$ e $F'(p, \alpha)$ dois símbolos de função binários introduzidos por definição e $F$ a $\alpha$-ésima iteração de $G$.
    Seja $p$ tal que $\forall \alpha \in \ON \Bigl( F'_p(\alpha)=G(p, F'_p  |\alpha)\Bigr)$.

    Então $\forall \alpha \in \ON \Bigl( F_p(\alpha)=F'_p(\alpha)\Bigr)$.
\end{theorem}
\begin{proof}
    Procedemos por indução transfinita sobre $\alpha$.
    Se o resultado vale para todo $\beta <\alpha$, então $F'_p|\alpha=F_p|\alpha$.
    Assim, $F'_p(\alpha)=G(p, F'_p|\alpha)=G(p, F_p|\alpha)=F_p(\alpha)$.
\end{proof}
